\chapter{Learning With Errors (LWE)}
\label{ch:lwe}

\section{From exact linear algebra to a noisy problem}

Solving a system of linear equations over a field is one of the oldest solved problems in mathematics. Given a matrix $\mathbf{A}\in\mathbb{Z}_q^{m\times n}$ and a vector $\mathbf{b}\in\mathbb{Z}_q^{m}$, Gaussian elimination recovers any $\mathbf{s}$ with $\mathbf{A}\mathbf{s}=\mathbf{b}$ in $O(n^3)$ field operations. Because the problem is this easy, an \emph{exact} linear system carries no cryptographic hardness whatsoever: whatever the secret $\mathbf{s}$ encodes, an adversary reads it off directly.

The Learning With Errors (LWE) problem, introduced by Regev~\cite{regev2009}, makes a single, deliberate change to this picture: it perturbs each equation by a small additive error. The observed system is
\[
\mathbf{b} = \mathbf{A}\mathbf{s} + \mathbf{e} \pmod q,
\]
where $\mathbf{e}$ is a short noise vector drawn from a prescribed distribution. Gaussian elimination is now useless -- it amplifies the noise catastrophically -- and the problem of recovering $\mathbf{s}$ becomes, as far as is known, exponentially hard. This chapter develops LWE as a formal object: its distribution, its two decision problems, the geometry that explains its hardness, and the worst-case reduction that makes it a trustworthy foundation. It is the assumption on which ML-KEM (Chapter~\ref{ch:fips203}) and ML-DSA (Chapter~\ref{ch:mldsa}) ultimately rest.

\section{The LWE distribution}

Throughout, fix a dimension $n$ (the security parameter), a modulus $q=q(n)$, and an \emph{error distribution} $\chi$ over $\mathbb{Z}_q$ concentrated on values of small absolute value. Vectors are column vectors, and $\langle\cdot,\cdot\rangle$ denotes the inner product modulo $q$.

\begin{definition}[LWE distribution]
For a fixed secret $\mathbf{s}\in\mathbb{Z}_q^{n}$, the \emph{LWE distribution} $A_{\mathbf{s},\chi}$ over $\mathbb{Z}_q^{n}\times\mathbb{Z}_q$ is sampled by drawing $\mathbf{a}\leftarrow\mathbb{Z}_q^{n}$ uniformly and $e\leftarrow\chi$, and outputting
\[
(\mathbf{a},\,b), \qquad b=\langle \mathbf{a},\mathbf{s}\rangle + e \pmod q.
\]
\end{definition}

Collecting $m$ independent samples and stacking the vectors $\mathbf{a}_i$ as the rows of a matrix $\mathbf{A}\in\mathbb{Z}_q^{m\times n}$ gives the compact matrix form
\[
(\mathbf{A},\,\mathbf{b}), \qquad \mathbf{b}=\mathbf{A}\mathbf{s}+\mathbf{e}\pmod q,\quad \mathbf{e}\leftarrow\chi^{m}.
\]
The name is now transparent: the adversary is asked to \emph{learn} the hidden $\mathbf{s}$ from many noisy linear measurements of it. The only obstruction is the error $\mathbf{e}$; without it the samples are an exactly solvable system.

\begin{figure}[H]
\centering
\begin{tikzpicture}[node distance=8mm and 16mm]
  \node[flownode] (A) {$\mathbf{A}$};
  \node[flownode, below=8mm of A] (s) {$\mathbf{s}$};
  \node[flownode, right=16mm of A, yshift=-4mm] (mul) {$\times$};
  \node[flownode, right=16mm of mul] (As) {$\mathbf{A}\mathbf{s}$};
  \node[flownode, below=8mm of As] (e) {$\mathbf{e}$};
  \node[flownode, right=16mm of As, yshift=-4mm] (add) {$+$};
  \node[keynode, right=16mm of add] (b) {$\mathbf{b}=\mathbf{A}\mathbf{s}+\mathbf{e}$};
  \draw[flowarrow] (A) -- (mul);
  \draw[flowarrow] (s) -- (mul);
  \draw[flowarrow] (mul) -- (As);
  \draw[flowarrow] (As) -- (add);
  \draw[flowarrow] (e) -- (add);
  \draw[flowarrow] (add) -- (b);
\end{tikzpicture}
\caption{The LWE sample equation as a computation graph: the public matrix $\mathbf{A}$ and the secret $\mathbf{s}$ feed a multiply node, and the noise $\mathbf{e}$ is added to the product $\mathbf{A}\mathbf{s}$ to produce the public value $\mathbf{b}$. All arithmetic is modulo $q$. Recovering $\mathbf{s}$ from $(\mathbf{A},\mathbf{b})$ alone is the LWE problem; the entire difficulty is contributed by the otherwise-negligible term $\mathbf{e}$.}
\label{fig:lwe-equation-boxes}
\end{figure}

\begin{remark}[On the secret distribution]
The definition draws $\mathbf{s}$ uniformly from $\mathbb{Z}_q^{n}$, but a standard transformation puts LWE into \emph{normal form}, in which the secret is drawn from the error distribution $\chi$ and is therefore itself short, without loss of hardness. Deployed schemes exploit this: in ML-KEM the secret is a small vector sampled from a centered binomial distribution (Chapter~\ref{ch:cbd}), which shortens keys and speeds arithmetic while preserving security.
\end{remark}

\section{The two LWE problems}

LWE gives rise to a search problem and a decision problem, and the distinction is central: encryption schemes are proved secure against the \emph{decision} version, while the \emph{search} version is the more intuitive statement of ``recover the key.''

\begin{definition}[Search-LWE]
Given $m$ independent samples $(\mathbf{A},\mathbf{b})$ drawn from $A_{\mathbf{s},\chi}$ for a fixed but unknown $\mathbf{s}\in\mathbb{Z}_q^{n}$, recover $\mathbf{s}$.
\end{definition}

\begin{definition}[Decision-LWE]
Distinguish, with non-negligible advantage, between the two distributions on $\mathbb{Z}_q^{m\times n}\times\mathbb{Z}_q^{m}$:
\begin{itemize}
  \item the LWE distribution: $(\mathbf{A},\mathbf{A}\mathbf{s}+\mathbf{e})$ with $\mathbf{s}\leftarrow\mathbb{Z}_q^{n}$ and $\mathbf{e}\leftarrow\chi^{m}$;
  \item the uniform distribution: $(\mathbf{A},\mathbf{u})$ with $\mathbf{u}\leftarrow\mathbb{Z}_q^{m}$ uniform and independent of $\mathbf{A}$.
\end{itemize}
\end{definition}

The decision formulation is what makes LWE a workhorse for cryptography: if LWE samples are computationally indistinguishable from uniform randomness, then a ciphertext that masks a message by adding it to an LWE sample is indistinguishable from a random string, which is exactly semantic security. Remarkably, the two problems are equivalent.

\begin{proposition}[Search--decision equivalence]
\label{prop:search-decision}
For a prime modulus $q=\mathrm{poly}(n)$, Search-LWE and Decision-LWE are equivalent up to a polynomial factor in running time and sample complexity: an efficient algorithm for either yields an efficient algorithm for the other.
\end{proposition}

The forward direction (search solves decision) is immediate. The reverse -- bootstrapping a mere distinguisher into a full secret-recovery algorithm -- is the substantial part: one guesses each coordinate $s_j\in\mathbb{Z}_q$ in turn, uses the distinguisher to test the guess, and amplifies. Regev proved the equivalence for prime $q=\mathrm{poly}(n)$~\cite{regev2009}; later work extended it to composite and larger moduli. The practical consequence is that designers may reason about the clean search problem while security proofs consume the decision problem.

\section{The error distribution and the modulus-to-noise ratio}

The error distribution $\chi$ is the beating heart of LWE. In Regev's formulation it is a \emph{discrete Gaussian} $D_{\mathbb{Z},\sigma}$ of width $\sigma$: the probability of an integer $x$ is proportional to $\exp(-\pi x^2/\sigma^2)$, rounded to $\mathbb{Z}$. It is convenient to parameterize the width relative to the modulus by a factor $\alpha\in(0,1)$,
\[
\sigma \;=\; \alpha q,
\]
so that $\alpha$ measures the noise as a fraction of the modulus. Two competing constraints pin down the right regime.

\begin{itemize}
  \item \textbf{Correctness needs small noise.} A legitimate decryptor removes the error by rounding, which succeeds only while $|e|$ stays below roughly $q/4$. Too much noise and honest decryption fails.
  \item \textbf{Security needs enough noise.} If $\alpha q$ is too small the error is negligible and the system is essentially exact, hence easy. The worst-case hardness theorem below requires $\alpha q \gtrsim \sqrt{n}$.
\end{itemize}

The single quantity that governs both is the \emph{modulus-to-noise ratio} $q/\sigma = 1/\alpha$. A larger ratio makes decryption more robust but the underlying problem easier (it corresponds to a larger lattice approximation factor); a smaller ratio does the reverse. Parameter selection for a real scheme is, at bottom, the art of choosing this ratio. In practice the exact discrete Gaussian is replaced by the centered binomial distribution (Chapter~\ref{ch:cbd}), which has essentially the same shape but is far easier to sample in constant time; the foundational analysis, however, is cleanest for the Gaussian~\cite{micciancioregev2007}.

\section{A geometric view: bounded-distance decoding}

LWE has a precise lattice interpretation, which is the source of both its hardness and its name-recognition among lattice problems. Given $\mathbf{A}\in\mathbb{Z}_q^{m\times n}$, define the \emph{$q$-ary lattice}
\[
\Lambda_q(\mathbf{A}) \;=\; \{\, \mathbf{A}\mathbf{s} \bmod q : \mathbf{s}\in\mathbb{Z}_q^{n} \,\} + q\mathbb{Z}^{m} \;\subseteq\; \mathbb{Z}^{m}.
\]
The noiseless product $\mathbf{A}\mathbf{s}$ is exactly a point of this lattice, and the observed vector $\mathbf{b}=\mathbf{A}\mathbf{s}+\mathbf{e}$ is that lattice point displaced by the short offset $\mathbf{e}$. Recovering $\mathbf{s}$ is therefore the problem of finding the lattice point nearest to $\mathbf{b}$.

This is the closest vector problem (CVP) of Chapter~\ref{ch:svp}, but with a crucial promise: the target $\mathbf{b}$ is guaranteed to lie \emph{within a small, known distance} of the lattice, since $\|\mathbf{e}\|$ is bounded by the error distribution. A CVP instance carrying this promise is called \emph{Bounded-Distance Decoding} (BDD).

\begin{proposition}[LWE is bounded-distance decoding]
Search-LWE with error norm bounded by $r$ is an instance of $\mathrm{BDD}$ on the lattice $\Lambda_q(\mathbf{A})$ with decoding radius $r$; the secret $\mathbf{s}$ is recovered from the closest lattice point $\mathbf{A}\mathbf{s}$.
\end{proposition}

The BDD promise is what an attacker tries to exploit and what the parameters are tuned to deny. When the decoding radius is a small fraction of the lattice's minimum distance, the nearest lattice point is unique but still hard to find in high dimension; this is the operating point of every deployed scheme.

\section{Hardness: the worst-case to average-case reduction}

The reason LWE is trusted is not merely that no efficient attack is known. It is that Regev proved a \emph{worst-case to average-case reduction}: breaking LWE on random instances is at least as hard as solving standard lattice problems on \emph{every} instance.

\begin{theorem}[Worst-case hardness of LWE, informal~\cite{regev2009}]
\label{thm:regev}
Let $n$ be a security parameter, $q=q(n)\le\mathrm{poly}(n)$ a modulus, and $\alpha=\alpha(n)\in(0,1)$ with $\alpha q>2\sqrt{n}$. If there is an efficient (possibly randomized) algorithm that solves Decision-LWE$_{n,q,\alpha}$ on the average, then there is an efficient \emph{quantum} algorithm that approximates the decision shortest-vector problem $(\mathrm{GapSVP})$ and the shortest-independent-vectors problem $(\mathrm{SIVP})$ to within $\widetilde{O}(n/\alpha)$ factors, in the worst case, on every $n$-dimensional lattice.
\end{theorem}

Three points give this theorem its weight.

\begin{itemize}
  \item \textbf{Worst-case to average-case.} A scheme's keys are random (average-case) LWE instances. The reduction guarantees that breaking a random key is as hard as the worst conceivable lattice, so there are no ``weak keys'' to stumble into -- a guarantee that number-theoretic assumptions such as factoring do not provide.
  \item \textbf{The role of quantumness.} The reduction uses a quantum algorithm as a proof device only. LWE-based schemes run on entirely classical hardware; the quantum step is in the security argument, not the cryptosystem. Peikert later gave a \emph{classical} reduction for exponential moduli~\cite{peikert2009}, and Brakerski, Langlois, Peikert, Regev, and Stehl\'{e} established classical hardness for polynomial moduli~\cite{blprs2013}, removing the quantum caveat.
  \item \textbf{The approximation factor.} The reduction yields hardness for $\widetilde{O}(n/\alpha)$-approximate lattice problems, which are believed to require time $2^{\Theta(n)}$ for the parameters used in practice. The factor $n/\alpha$ makes precise the tension of the previous section: less noise (smaller $\alpha$) yields a weaker guarantee.
\end{itemize}

\section{The dual problem: Short Integer Solution (SIS)}
\label{sec:sis}

LWE has a close relative that becomes important once we reach signatures. Where LWE \emph{hides} a secret inside noise and asks you to recover it, the \textbf{Short Integer Solution} (SIS) problem~\cite{ajtai1996} asks you to \emph{find} a short combination of public vectors. Given a random matrix
\[
\mathbf{A} \in \mathbb{Z}_q^{n \times m}, \qquad m > n,
\]
the task is to find a short nonzero integer vector $\mathbf{z}$ with
\[
\mathbf{A}\mathbf{z} = \mathbf{0} \pmod q, \qquad 0 < \|\mathbf{z}\| \le \beta.
\]
Because there are more columns than rows, such vectors certainly exist; the difficulty is finding one that is \emph{short}. A random solution would have large entries, and shrinking it is exactly the kind of short-vector search that is hard in a high-dimensional lattice. Geometrically, SIS lives in the \emph{orthogonal} $q$-ary lattice
\[
\Lambda_q^{\perp}(\mathbf{A}) \;=\; \{\, \mathbf{z}\in\mathbb{Z}^{m} : \mathbf{A}\mathbf{z}=\mathbf{0}\bmod q \,\},
\]
and a solution is a short vector of $\Lambda_q^{\perp}(\mathbf{A})$ -- the SVP counterpart to the BDD instance $\Lambda_q(\mathbf{A})$ that defines LWE.

The two problems are mirror images. LWE is a \emph{search among noisy points} (recover $\mathbf{s}$ from $\mathbf{A}\mathbf{s}+\mathbf{e}$); SIS is a \emph{search for a short null vector} (find $\mathbf{z}$ with $\mathbf{A}\mathbf{z}=\mathbf{0}$). Both are hard for the same underlying reason -- short vectors in the associated lattice are hard to find -- and both admit worst-case reductions (SIS to the worst case of SIVP~\cite{ajtai1996}). We will meet SIS again directly: in ML-DSA / Dilithium (Chapter~\ref{ch:mldsa}), forging a signature amounts to producing a short vector satisfying a public equation, which is precisely an SIS instance.

\section{Concrete security and cryptanalysis}

Theorem~\ref{thm:regev} is asymptotic; deploying LWE requires \emph{concrete} parameters $(n,q,\alpha,m)$ for which the best known attack costs more than $2^{128}$ (or $2^{192}$, $2^{256}$) operations. Concrete security is estimated against lattice-reduction attacks rather than the reduction itself, since the reduction is not tight.

\begin{itemize}
  \item \textbf{Primal attack.} Embed the BDD instance of $\S$5 as a unique-shortest-vector instance in a lattice of dimension $\approx m+n+1$ (Kannan's embedding) and run lattice reduction to find the short error vector, revealing $\mathbf{s}$.
  \item \textbf{Dual attack.} Find short vectors in the orthogonal lattice $\Lambda_q^{\perp}(\mathbf{A})$ and use them to distinguish LWE samples from uniform, solving Decision-LWE directly.
\end{itemize}

Both are driven by the \textbf{Block-Korkine-Zolotarev} (BKZ) algorithm with block size $\beta$, whose cost grows as $2^{\Theta(\beta)}$; the attacker's task reduces to finding the smallest $\beta$ that achieves the root-Hermite factor needed to solve the instance. The community consolidates these estimates in the ``LWE estimator''~\cite{aps2015}, building on the concrete-attack analysis of Lindner and Peikert~\cite{lindnerpeikert2011}; ML-KEM's parameter sets are chosen to clear the target security levels under this estimator.

\begin{remark}[The number of samples is part of the assumption]
Attack cost depends on how many samples $m$ the adversary sees, not only on $(n,q,\alpha)$: more equations can make the underlying lattice easier to reduce. A well-designed scheme therefore \emph{bounds} the number of LWE samples an adversary can obtain from public data. This is one reason a KEM publishes a fixed-size key rather than an unbounded stream of samples.
\end{remark}

\section{Choosing the modulus $q$}

The modulus ties the preceding threads together. It must be large enough that the error rarely wraps around modulo $q$ and destroys correctness, yet small enough to keep keys compact and arithmetic fast; and, as $\S$4 showed, the ratio $q/\sigma$ directly sets the security--correctness trade-off. Real schemes add a further constraint: $q$ is chosen to admit an efficient number-theoretic transform (Chapter~\ref{ch:ntt}), which is why ML-KEM fixes the specific prime $q=3329$. Modulus selection is thus the first place where the abstract assumption meets hard engineering.

\section{Sage experiments}

The quickest route to intuition is to build a small instance by hand and watch the error convert an easy system into a hard one. Fix a modulus and dimension.

\begin{lstlisting}[language=Python]
q = 17
n = 3
F = GF(q)
\end{lstlisting}

Choose a secret vector and a batch of random rows.

\begin{lstlisting}[language=Python]
from sage.all import *
s = random_vector(F, n)
A = random_matrix(F, 8, n)
print("secret =", s)
\end{lstlisting}

Without noise, $\mathbf{b}=\mathbf{A}\mathbf{s}$ is an exact system and $\mathbf{s}$ is recovered by solving it directly.

\begin{lstlisting}[language=Python]
b = A * s
print(A.solve_right(b) == s)   # True: exact recovery
\end{lstlisting}

Now add a short error. In $\mathbb{F}_{17}$ the value $16$ represents $-1$, so this vector has entries in $\{-1,0,1\}$.

\begin{lstlisting}[language=Python]
noise   = vector(F, [1, 0, 16, 1, 0, 0, 16, 1])
b_noisy = A * s + noise
\end{lstlisting}

The system $\mathbf{A}\mathbf{s}'=\mathbf{b}_{\text{noisy}}$ now has no exact solution, and naive solving returns the wrong secret. Sweeping the error magnitude reveals the phase transition of $\S$4: below a noise threshold the secret is still recoverable by more careful decoding, and above it recovery becomes infeasible. Algorithm~\ref{alg:lwe-sample} states the sampling procedure formally, using the same variable names as the code.

\begin{algorithm}[H]
\caption{Generate an LWE Sample}
\label{alg:lwe-sample}
\begin{algorithmic}[1]
\Require Modulus $q$; dimension $n$; number of rows $m$; error distribution $\chi$ over $\mathbb{Z}_q$
\Ensure Public sample $(\mathbf{A},\mathbf{b})$ and secret $\mathbf{s}$
\State $\mathbf{s} \gets$ uniformly random vector in $\mathbb{Z}_q^{n}$ \Comment{\texttt{s = random\_vector(F, n)}}
\State $\mathbf{A} \gets$ uniformly random matrix in $\mathbb{Z}_q^{m \times n}$ \Comment{\texttt{A = random\_matrix(F, m, n)}}
\State $\mathbf{e} \gets (e_1,\dots,e_m)$ with each $e_i \gets \chi$ \Comment{\texttt{noise = vector(F, [...])}}
\State $\mathbf{b} \gets \mathbf{A}\mathbf{s} + \mathbf{e} \pmod q$ \Comment{\texttt{b\_noisy = A * s + noise}}
\State \Return public sample $(\mathbf{A},\mathbf{b})$, secret $\mathbf{s}$
\end{algorithmic}
\end{algorithm}

\section{Summary}

LWE turns the easiest problem in linear algebra into one of the most robust hardness assumptions in cryptography by adding calibrated noise:
\[
\mathbf{b}=\mathbf{A}\mathbf{s}+\mathbf{e}\pmod q.
\]
Its structure rests on a few precise ideas: the search and decision formulations are equivalent (Proposition~\ref{prop:search-decision}); the problem is bounded-distance decoding on the $q$-ary lattice $\Lambda_q(\mathbf{A})$; its average-case hardness follows from worst-case lattice problems (Theorem~\ref{thm:regev}); its dual SIS lives in the orthogonal lattice $\Lambda_q^{\perp}(\mathbf{A})$; and concrete parameters are set against BKZ-based estimators rather than the asymptotic reduction. The conceptual chain built over the last several chapters is now complete.

\begin{center}
\begin{tikzpicture}[node distance=5mm and 5mm, every node/.style={font=\scriptsize}]
  \node[flownode] (a) {finite field};
  \node[flownode, right=of a] (b) {linear\\[-2pt]algebra};
  \node[flownode, right=of b] (c) {lattice};
  \node[flownode, right=of c] (d) {SVP / CVP};
  \node[flownode, right=of d] (e) {noise};
  \node[keynode, right=of e] (f) {LWE};
  \draw[flowarrow] (a) -- (b);
  \draw[flowarrow] (b) -- (c);
  \draw[flowarrow] (c) -- (d);
  \draw[flowarrow] (d) -- (e);
  \draw[flowarrow] (e) -- (f);
\end{tikzpicture}
\end{center}

The next chapter carries this assumption into a ring, replacing vectors of integers with polynomials to obtain the far more efficient Ring-LWE.
